% LaTeX Template for AI Problem Analysis Output
% AMC12A 2024 Problem 10 - Strategic Analysis

\documentclass[11pt]{article}
\usepackage[utf8]{inputenc}
\usepackage{amsmath, amssymb, amsthm}
\usepackage{geometry}
\usepackage{xcolor}
\usepackage[most]{tcolorbox}
\usepackage{enumitem}
\usepackage{fancyhdr}

% Page setup
\geometry{margin=1in}
\setlength{\headheight}{14.5pt}
\pagestyle{fancy}
\fancyhf{}
\rhead{AMC12/COMC Problem Analysis}
\lhead{2024 AMC 12A Problem 10}
\cfoot{\thepage}

% Color definitions
\definecolor{problemconfig}{RGB}{230, 242, 255}    % Light blue
\definecolor{strategic}{RGB}{255, 230, 242}       % Light pink
\definecolor{tactical}{RGB}{242, 255, 230}        % Light green
\definecolor{techniques}{RGB}{255, 242, 230}      % Light yellow
\definecolor{verification}{RGB}{255, 230, 230}    % Light red
\definecolor{solution}{RGB}{230, 255, 255}        % Light cyan
\definecolor{competition}{RGB}{242, 242, 255}     % Light purple
\definecolor{gold}{RGB}{218, 165, 32}

% Box styles
\newtcolorbox{problemconfigbox}{
    colback=problemconfig,
    colframe=blue,
    title=PROBLEM CONFIGURATION ANALYSIS,
    fonttitle=\bfseries,
    arc=3pt,
    boxrule=1pt,
    breakable
}

\newtcolorbox{strategicbox}{
    colback=strategic,
    colframe=purple,
    title=STRATEGIC FRAMEWORK APPLICATION,
    fonttitle=\bfseries,
    arc=3pt,
    boxrule=1pt,
    breakable
}

\newtcolorbox{tacticalbox}{
    colback=tactical,
    colframe=green,
    title=TACTICAL METHODOLOGY SELECTION,
    fonttitle=\bfseries,
    arc=3pt,
    boxrule=1pt,
    breakable
}

\newtcolorbox{techniquesbox}{
    colback=techniques,
    colframe=orange,
    title=CONVERSION TECHNIQUES APPLICATION,
    fonttitle=\bfseries,
    arc=3pt,
    boxrule=1pt,
    breakable
}

\newtcolorbox{verificationbox}{
    colback=verification,
    colframe=red,
    title=VERIFICATION STRATEGY DESIGN,
    fonttitle=\bfseries,
    arc=3pt,
    boxrule=1pt,
    breakable
}

\newtcolorbox{solutionbox}{
    colback=solution,
    colframe=cyan,
    title=SOLUTION ROADMAP,
    fonttitle=\bfseries,
    arc=3pt,
    boxrule=1pt,
    breakable
}

\newtcolorbox{competitionbox}{
    colback=competition,
    colframe=purple,
    title=COMPETITION STRATEGY INTEGRATION,
    fonttitle=\bfseries,
    arc=3pt,
    boxrule=1pt,
    breakable
}

\newtcolorbox{keyinsightbox}{
    colback=yellow!20,
    colframe=gold,
    title=KEY INSIGHT,
    fonttitle=\bfseries,
    arc=3pt,
    boxrule=2pt,
    leftrule=5pt,
    breakable
}

\newtcolorbox{warningbox}{
    colback=red!10,
    colframe=red!50,
    title=WARNING: Common Pitfall,
    fonttitle=\bfseries,
    arc=3pt,
    boxrule=2pt,
    leftrule=5pt,
    breakable
}

\newtcolorbox{tipbox}{
    colback=green!10,
    colframe=green!50,
    title=PRO TIP,
    fonttitle=\bfseries,
    arc=3pt,
    boxrule=2pt,
    leftrule=5pt,
    breakable
}

\begin{document}

\title{\textbf{AMC12A 2024 Problem 10\\Strategic Analysis}}
\author{Competition Math Framework v2.1}
\date{\today}
\maketitle

\section*{Problem Statement}

Let $\alpha$ be the radian measure of the smallest angle in a $3{-}4{-}5$ right triangle. Let $\beta$ be the radian measure of the smallest angle in a $7{-}24{-}25$ right triangle. In terms of $\alpha$, what is $\beta$?

\vspace{0.5em}
\textbf{Answer Choices:}\\
$\textbf{(A) }\frac{\alpha}{3}\qquad \textbf{(B) }\alpha - \frac{\pi}{8}\qquad \textbf{(C) }\frac{\pi}{2} - 2\alpha \qquad \textbf{(D) }\frac{\alpha}{2}\qquad \textbf{(E) }\pi - 4\alpha$

\begin{problemconfigbox}
\subsection*{A. Problem Classification}
\begin{itemize}[leftmargin=*]
    \item \textbf{Problem Type}: To Find - express one angle in terms of another
    \item \textbf{Competition Context}: AMC12A (75 min, 25 problems) - Problem 10 of 25
    \item \textbf{Difficulty Band}: Mid-band (P9-16) - requires trigonometric identities and pattern recognition
    \item \textbf{Mathematical Domain}: Trigonometry with Geometry foundation (right triangles)
    \item \textbf{Estimated Time}: 2-3 minutes for students familiar with trig identities; 4-5 minutes otherwise
\end{itemize}

\subsection*{B. Problem Structure Identification}
\begin{itemize}[leftmargin=*]
    \item \textbf{Given Information}: 
    \begin{itemize}
        \item Triangle 1: $3{-}4{-}5$ right triangle (Pythagorean triple: $3^2 + 4^2 = 9 + 16 = 25 = 5^2$)
        \item Triangle 2: $7{-}24{-}25$ right triangle (Pythagorean triple: $7^2 + 24^2 = 49 + 576 = 625 = 25^2$)
        \item $\alpha$ = smallest angle in triangle 1
        \item $\beta$ = smallest angle in triangle 2
        \item Both angles measured in radians
    \end{itemize}
    
    \item \textbf{Constraints}: 
    \begin{itemize}
        \item $\alpha$ is opposite the shortest side (3) in the $3{-}4{-}5$ triangle
        \item $\beta$ is opposite the shortest side (7) in the $7{-}24{-}25$ triangle
        \item Both triangles are right triangles
        \item Answer must be expressed in terms of $\alpha$
    \end{itemize}
    
    \item \textbf{Objective}: Express $\beta$ as a function of $\alpha$ using the given answer choices
    
    \item \textbf{Hidden Assumptions}: 
    \begin{itemize}
        \item The smallest angle is opposite the smallest side (standard triangle properties)
        \item Radian measure implies use of trigonometric functions without degree conversion
        \item The relationship involves standard trig identities or angle sum formulas
        \item Answer choices suggest relationships involving $\pi/2$ (complementary angles) and multiples of $\alpha$
    \end{itemize}
    
    \item \textbf{Answer Choice Leverage}: 
    \begin{itemize}
        \item Choice (C) $\frac{\pi}{2} - 2\alpha$ involves complementary angle and double angle - suggests cofunction identity
        \item Choice (D) $\frac{\alpha}{2}$ is simple half-angle - easily testable
        \item Choice (A) $\frac{\alpha}{3}$ suggests specific ratio relationship
        \item Choices (B) and (E) involve specific fractions of $\pi$ - less natural
        \item Can test choices by computing $\sin$ or $\tan$ of both sides
    \end{itemize}
\end{itemize}

\subsection*{C. Strategic Orientation}
\begin{itemize}[leftmargin=*]
    \item \textbf{Hypothesis}: Given two Pythagorean triple right triangles with angles $\alpha$ and $\beta$
    \item \textbf{Conclusion}: Determine the relationship $\beta = f(\alpha)$
    \item \textbf{Notation}: 
    \begin{itemize}
        \item For $3{-}4{-}5$ triangle: $\sin\alpha = \frac{3}{5}$, $\cos\alpha = \frac{4}{5}$, $\tan\alpha = \frac{3}{4}$
        \item For $7{-}24{-}25$ triangle: $\sin\beta = \frac{7}{25}$, $\cos\beta = \frac{24}{25}$, $\tan\beta = \frac{7}{24}$
        \item Complementary angles: $\frac{\pi}{2} - \alpha$ is the larger acute angle in the $3{-}4{-}5$ triangle
    \end{itemize}
    
    \item \textbf{Intuitive Approach}: 
    \begin{itemize}
        \item Compute trig values for both angles
        \item Look for relationship using double angle formulas: $\sin(2\alpha) = 2\sin\alpha\cos\alpha$, $\cos(2\alpha) = \cos^2\alpha - \sin^2\alpha$
        \item Check if $\beta$ relates to $2\alpha$ or $\alpha/2$
        \item Use complementary angle: if $\cos(2\alpha) = \sin\beta$, then $2\alpha = \frac{\pi}{2} - \beta$
    \end{itemize}
    
    \item \textbf{Cross-Domain Potential}: 
    \begin{itemize}
        \item \textbf{Pure Trigonometry}: Use trig identities and formulas
        \item \textbf{Algebraic}: Test answer choices by computing and comparing values
        \item \textbf{Geometric}: Combine triangles physically to find angle relationships
        \item \textbf{Complex Numbers}: Represent angles as arguments of complex numbers
    \end{itemize}
\end{itemize}
\end{problemconfigbox}

\begin{strategicbox}
\subsection*{A. Psychological Strategies}
\begin{itemize}[leftmargin=*]
    \item \textbf{Mental Toughness}: This is a mid-band problem that appears straightforward but requires recognizing the right identity. Don't panic if the first approach doesn't yield immediate results.
    
    \item \textbf{Creativity Cultivation}: Consider multiple representations of angles - not just direct computation but also complementary angles, double angles, and half angles.
    
    \item \textbf{Iterative Refinement}: If direct computation of $\beta$ in terms of $\alpha$ is unclear, work backwards from answer choices or look for patterns in the trig values.
\end{itemize}

\subsection*{B. Investigation Strategies}
\begin{itemize}[leftmargin=*]
    \item \textbf{Startup Orientation}: 
    \begin{itemize}
        \item Immediately compute basic trig ratios: $\sin\alpha = 3/5$, $\cos\alpha = 4/5$, $\sin\beta = 7/25$, $\cos\beta = 24/25$
        \item Notice that $7/25 = 0.28$ and $24/25 = 0.96$ - these values should trigger pattern recognition
    \end{itemize}
    
    \item \textbf{Get Your Hands Dirty}: 
    \begin{itemize}
        \item Compute $\sin(2\alpha) = 2 \cdot \frac{3}{5} \cdot \frac{4}{5} = \frac{24}{25}$ - this equals $\cos\beta$!
        \item Compute $\cos(2\alpha) = \left(\frac{4}{5}\right)^2 - \left(\frac{3}{5}\right)^2 = \frac{16}{25} - \frac{9}{25} = \frac{7}{25}$ - this equals $\sin\beta$!
        \item This strongly suggests $\beta$ and $2\alpha$ are complementary angles
    \end{itemize}
    
    \item \textbf{Penultimate Step}: Before expressing $\beta$ in terms of $\alpha$, recognize that:
    \begin{itemize}
        \item If $\cos(2\alpha) = \sin\beta$, then by cofunction identity: $\cos(2\alpha) = \sin\left(\frac{\pi}{2} - 2\alpha\right)$
        \item Therefore $\sin\beta = \sin\left(\frac{\pi}{2} - 2\alpha\right)$
        \item Since both angles are in $(0, \frac{\pi}{2})$, we have $\beta = \frac{\pi}{2} - 2\alpha$
    \end{itemize}
    
    \item \textbf{Wishful Thinking}: ``I wish I could relate $\beta$ directly to $\alpha$... What if I try computing $2\alpha$?''
    
    \item \textbf{Conjecture via Small Cases}: 
    \begin{itemize}
        \item Check if $\beta = \alpha/2$: $\sin(\alpha/2) = \sqrt{\frac{1-\cos\alpha}{2}} = \sqrt{\frac{1-4/5}{2}} = \sqrt{\frac{1/5}{2}} = \sqrt{\frac{1}{10}} \neq \frac{7}{25}$ \texttimes
        \item Try $\beta = \frac{\pi}{2} - 2\alpha$: Use cofunction identity to verify \checkmark
    \end{itemize}
    
    \item \textbf{Visualization}: 
    \begin{itemize}
        \item Can geometrically construct both triangles and look for angle relationships
        \item Scale the $3{-}4{-}5$ triangle by factor of 6 to get $18{-}24{-}30$ triangle
        \item Combine with $7{-}24{-}25$ triangle sharing the side of length 24
        \item This creates isosceles triangle revealing angle relationships
    \end{itemize}
\end{itemize}

\subsection*{C. Argument Construction}
\begin{itemize}[leftmargin=*]
    \item \textbf{Proof Method}: Direct computation using double angle formulas and cofunction identities
    
    \item \textbf{Key Lemmas}: 
    \begin{itemize}
        \item \textbf{Double Angle Formulas}: $\cos(2\alpha) = \cos^2\alpha - \sin^2\alpha = 2\cos^2\alpha - 1 = 1 - 2\sin^2\alpha$
        \item \textbf{Cofunction Identity}: $\cos\theta = \sin\left(\frac{\pi}{2} - \theta\right)$ for all $\theta$
        \item \textbf{Angle Sum Formula}: $\tan(\alpha + \beta) = \frac{\tan\alpha + \tan\beta}{1 - \tan\alpha\tan\beta}$
    \end{itemize}
    
    \item \textbf{Logical Structure}: 
    \begin{enumerate}
        \item Compute $\tan\alpha = 3/4$ and $\tan\beta = 7/24$
        \item Compute $\tan(\alpha + \beta)$ using angle sum formula
        \item Show that $\tan(\alpha + \beta) = 4/3$
        \item Recognize that $\tan(\alpha + \beta) = 4/3 = \cot\alpha = \tan\left(\frac{\pi}{2} - \alpha\right)$
        \item Conclude that $\alpha + \beta = \frac{\pi}{2} - \alpha$
        \item Solve for $\beta$: $\beta = \frac{\pi}{2} - 2\alpha$
    \end{enumerate}
\end{itemize}
\end{strategicbox}

\begin{tacticalbox}
\subsection*{A. Core Mathematical Tactics}
\begin{itemize}[leftmargin=*]
    \item \textbf{Pattern Recognition}: The specific values $\sin\beta = 7/25$ and $\cos(2\alpha) = 7/25$ matching is the key pattern
    
    \item \textbf{Identity Application}: Systematic use of trigonometric identities (double angle, cofunction)
    
    \item \textbf{Complementary Angle Exploitation}: Recognizing that complementary angles have swapped sine and cosine values
    
    \item \textbf{Answer Choice Testing}: For AMC multiple choice, testing answer choices can be faster than full derivation
\end{itemize}

\subsection*{B. Domain-Specific Tactics}
\begin{itemize}[leftmargin=*]
    \item \textbf{Trigonometry}: 
    \begin{itemize}
        \item Computing exact trig values from right triangles
        \item Applying double angle formulas systematically
        \item Using cofunction identities for complementary angles
        \item Recognizing Pythagorean triples yields exact rational trig values
    \end{itemize}
    
    \item \textbf{Geometry}: 
    \begin{itemize}
        \item Identifying smallest angles in right triangles
        \item Understanding angle relationships in combined triangles
        \item Using similarity and scaling to create geometric relationships
    \end{itemize}
    
    \item \textbf{Algebra}: 
    \begin{itemize}
        \item Manipulating expressions to match target forms
        \item Solving for unknown angles in terms of known angles
    \end{itemize}
\end{itemize}

\subsection*{C. Crossover Techniques}
\begin{itemize}[leftmargin=*]
    \item \textbf{Trigonometric-Algebraic}: Express trig functions algebraically and manipulate equations
    
    \item \textbf{Geometric-Trigonometric}: Use geometric construction of triangles to derive trig relationships
    
    \item \textbf{Numerical-Symbolic}: Compute numerical values to guide symbolic manipulation
\end{itemize}
\end{tacticalbox}

\begin{techniquesbox}
\subsection*{A. High-Probability Techniques}
\begin{itemize}[leftmargin=*]
    \item \textbf{Primary Technique - Double Angle Formula (H-Level)}: 
    \begin{itemize}
        \item \textbf{Recognition Cue}: Answer choice (C) involves $2\alpha$, suggesting double angle
        \item \textbf{Application}: Compute $\cos(2\alpha) = \cos^2\alpha - \sin^2\alpha = \frac{16}{25} - \frac{9}{25} = \frac{7}{25} = \sin\beta$
        \item \textbf{Success Rate}: 85-95\% for students who recognize double angle patterns
    \end{itemize}
    
    \item \textbf{Secondary Technique - Angle Addition Formula (H2)}: 
    \begin{itemize}
        \item \textbf{Recognition Cue}: Two angles given, need to find relationship
        \item \textbf{Application}: $\tan(\alpha + \beta) = \frac{3/4 + 7/24}{1 - (3/4)(7/24)} = \frac{18/24 + 7/24}{1 - 21/96} = \frac{25/24}{75/96} = \frac{25/24 \cdot 96/75} = 4/3$
        \item \textbf{Success Rate}: 70-80\% - more algebraically intensive
    \end{itemize}
    
    \item \textbf{Alternative - Answer Choice Testing (H17)}: 
    \begin{itemize}
        \item \textbf{Recognition Cue}: Multiple choice format with specific forms
        \item \textbf{Application}: Test (C) by computing $\sin\left(\frac{\pi}{2} - 2\alpha\right) = \cos(2\alpha) = 7/25 = \sin\beta$ \checkmark
        \item \textbf{Success Rate}: 90\%+ if done carefully
    \end{itemize}
\end{itemize}

\subsection*{B. Technique Selection Strategy}
\begin{itemize}[leftmargin=*]
    \item \textbf{Optimal Approach}: Use double angle formula to compute $\cos(2\alpha)$, recognize it equals $\sin\beta$, apply cofunction identity
    
    \item \textbf{Backup Approach}: Use angle addition formula for $\tan(\alpha + \beta)$, show $\alpha + \beta + \alpha = \pi/2$
    
    \item \textbf{Quick Check Approach}: Test answer choice (C) by computing both sides numerically or symbolically
    
    \item \textbf{Time Efficiency}: 
    \begin{itemize}
        \item Double angle method: 1-2 minutes
        \item Angle addition method: 2-3 minutes
        \item Answer testing: 1.5-2 minutes
    \end{itemize}
\end{itemize}

\subsection*{C. Technique Integration}
\begin{itemize}[leftmargin=*]
    \item \textbf{Integration Plan}:
    \begin{enumerate}
        \item Compute basic trig ratios for both angles (30 seconds)
        \item Try double angle formula on $\alpha$ (30 seconds)
        \item Recognize $\cos(2\alpha) = \sin\beta$ (10 seconds)
        \item Apply cofunction identity (20 seconds)
        \item Verify answer matches choice (C) (10 seconds)
    \end{enumerate}
    
    \item \textbf{Conflict Resolution}: If double angle doesn't immediately work, pivot to answer testing rather than trying all formulas
\end{itemize}
\end{techniquesbox}

\begin{verificationbox}
\subsection*{A. Verification Level Selection}

Using the 7-level verification system:

\begin{itemize}[leftmargin=*]
    \item \textbf{Selected Level}: Level 3 - Targeted Spot Check (15-30 seconds)
    \item \textbf{Rationale}: 
    \begin{itemize}
        \item Mid-band AMC problem (P10) with clear solution path
        \item Trig identity application is either right or wrong - intermediate verification less critical
        \item Can quickly verify by checking $\sin\beta = \sin\left(\frac{\pi}{2} - 2\alpha\right)$
        \item 15-30 seconds sufficient for confidence boost
    \end{itemize}
\end{itemize}

\subsection*{B. Verification Checklist}
\begin{itemize}[leftmargin=*]
    \item \textbf{Numerical Verification}: 
    \begin{itemize}
        \item Compute $\sin\beta = \frac{7}{25} = 0.28$
        \item Compute $\sin\left(\frac{\pi}{2} - 2\alpha\right) = \cos(2\alpha) = \frac{16-9}{25} = \frac{7}{25}$ \checkmark
        \item Values match exactly
    \end{itemize}
    
    \item \textbf{Alternative Check}: 
    \begin{itemize}
        \item Check angle sum: If $\beta = \frac{\pi}{2} - 2\alpha$, then $2\alpha + \beta = \frac{\pi}{2}$
        \item Verify: $\sin(2\alpha + \beta) = \sin(\pi/2) = 1$
        \item Using angle sum: $\sin(2\alpha)\cos\beta + \cos(2\alpha)\sin\beta = \frac{24}{25} \cdot \frac{24}{25} + \frac{7}{25} \cdot \frac{7}{25} = \frac{576 + 49}{625} = \frac{625}{625} = 1$ \checkmark
    \end{itemize}
    
    \item \textbf{Answer Choice Elimination}:
    \begin{itemize}
        \item (A) $\alpha/3$: $\sin(\alpha/3) \approx \sin(0.21) \approx 0.21 \neq 0.28$ \texttimes
        \item (D) $\alpha/2$: $\sin(\alpha/2) = \sqrt{(1-4/5)/2} = \sqrt{1/10} \approx 0.316 \neq 0.28$ \texttimes
        \item (E) $\pi - 4\alpha$: Would give $\beta > \pi/2$ for small $\alpha$, impossible for acute angle \texttimes
        \item (B) $\alpha - \pi/8$: $\alpha \approx 0.64$ rad, $\beta \approx 0.28$ rad, $\alpha - \pi/8 \approx 0.25 \neq 0.28$ \texttimes
        \item (C) $\frac{\pi}{2} - 2\alpha$: Verified above \checkmark
    \end{itemize}
\end{itemize}

\subsection*{C. Confidence Calibration}
\begin{itemize}[leftmargin=*]
    \item \textbf{Target Confidence}: 95\%+
    \item \textbf{Confidence Indicators}:
    \begin{itemize}
        \item[$\checkmark$] Double angle formula yielded exact match
        \item[$\checkmark$] Cofunction identity cleanly applied
        \item[$\checkmark$] Answer matches natural form in choice (C)
        \item[$\checkmark$] Alternative verification confirms result
        \item[$\checkmark$] All other answer choices eliminated
    \end{itemize}
    \item \textbf{Red Flags}: None - solution is clean and verifiable
\end{itemize}
\end{verificationbox}

\begin{solutionbox}
\subsection*{A. Step-by-Step Solution}

\textbf{Method 1: Double Angle Formula (Primary)}

\textbf{Step 1: Identify trigonometric ratios}

For the $3{-}4{-}5$ right triangle with smallest angle $\alpha$:
\begin{itemize}
    \item $\sin\alpha = \frac{3}{5}$ (opposite/hypotenuse)
    \item $\cos\alpha = \frac{4}{5}$ (adjacent/hypotenuse)
    \item $\tan\alpha = \frac{3}{4}$ (opposite/adjacent)
\end{itemize}

For the $7{-}24{-}25$ right triangle with smallest angle $\beta$:
\begin{itemize}
    \item $\sin\beta = \frac{7}{25}$ (opposite/hypotenuse)
    \item $\cos\beta = \frac{24}{25}$ (adjacent/hypotenuse)
    \item $\tan\beta = \frac{7}{24}$ (opposite/adjacent)
\end{itemize}

\textbf{Step 2: Apply double angle formula}

Compute $\cos(2\alpha)$ using the double angle formula:
\[
\cos(2\alpha) = \cos^2\alpha - \sin^2\alpha = \left(\frac{4}{5}\right)^2 - \left(\frac{3}{5}\right)^2 = \frac{16}{25} - \frac{9}{25} = \frac{7}{25}
\]

\textbf{Step 3: Recognize the match}

Observe that:
\[
\cos(2\alpha) = \frac{7}{25} = \sin\beta
\]

\textbf{Step 4: Apply cofunction identity}

By the cofunction identity, $\cos\theta = \sin\left(\frac{\pi}{2} - \theta\right)$ for any angle $\theta$.

Therefore:
\[
\cos(2\alpha) = \sin\left(\frac{\pi}{2} - 2\alpha\right)
\]

\textbf{Step 5: Conclude}

Since $\sin\beta = \cos(2\alpha) = \sin\left(\frac{\pi}{2} - 2\alpha\right)$ and both $\beta$ and $\frac{\pi}{2} - 2\alpha$ are angles in the interval $(0, \frac{\pi}{2})$ where sine is injective, we have:

\[
\beta = \frac{\pi}{2} - 2\alpha
\]

\textbf{Answer: (C) $\frac{\pi}{2} - 2\alpha$}

\subsection*{B. Alternative Approaches}

\textbf{Method 2: Angle Addition Formula}

\textbf{Step 1: Compute $\tan(\alpha + \beta)$}

Using the angle addition formula:
\[
\tan(\alpha + \beta) = \frac{\tan\alpha + \tan\beta}{1 - \tan\alpha\tan\beta} = \frac{\frac{3}{4} + \frac{7}{24}}{1 - \frac{3}{4} \cdot \frac{7}{24}}
\]

\textbf{Step 2: Simplify numerator}
\[
\frac{3}{4} + \frac{7}{24} = \frac{18}{24} + \frac{7}{24} = \frac{25}{24}
\]

\textbf{Step 3: Simplify denominator}
\[
1 - \frac{3}{4} \cdot \frac{7}{24} = 1 - \frac{21}{96} = 1 - \frac{7}{32} = \frac{32 - 7}{32} = \frac{25}{32}
\]

\textbf{Step 4: Compute ratio}
\[
\tan(\alpha + \beta) = \frac{25/24}{25/32} = \frac{25}{24} \cdot \frac{32}{25} = \frac{32}{24} = \frac{4}{3}
\]

\textbf{Step 5: Recognize complementary angle}

Note that $\tan\left(\frac{\pi}{2} - \alpha\right) = \cot\alpha = \frac{1}{\tan\alpha} = \frac{1}{3/4} = \frac{4}{3}$.

Therefore:
\[
\tan(\alpha + \beta) = \frac{4}{3} = \tan\left(\frac{\pi}{2} - \alpha\right)
\]

Since both angles are in $(0, \pi)$, we have:
\[
\alpha + \beta = \frac{\pi}{2} - \alpha
\]

Solving for $\beta$:
\[
\beta = \frac{\pi}{2} - 2\alpha
\]

\textbf{Method 3: Geometric Construction}

Scale the $3{-}4{-}5$ triangle by factor 6 to get $18{-}24{-}30$ triangle. Attach its side of length 24 to the corresponding side of the $7{-}24{-}25$ triangle. This creates a configuration where:
\begin{itemize}
    \item The combined base is $18 + 7 = 25$
    \item The height is 24
    \item The hypotenuse from far left to far right is 30 (scaled triangle) and 25 (other triangle)
    \item Triangle with base 25 and two equal sides (25 and 25) is isosceles
    \item Angle analysis shows $\alpha + \beta + \alpha = \frac{\pi}{2}$
    \item Therefore $\beta = \frac{\pi}{2} - 2\alpha$
\end{itemize}

\textbf{Method 4: Answer Choice Testing}

Test choice (C): $\beta = \frac{\pi}{2} - 2\alpha$

Verify: $\sin\beta = \sin\left(\frac{\pi}{2} - 2\alpha\right) = \cos(2\alpha) = \frac{16-9}{25} = \frac{7}{25}$ \checkmark

This matches $\sin\beta = \frac{7}{25}$ directly. $\checkmark$

\end{solutionbox}

\begin{competitionbox}
\subsection*{A. Time Management}
\begin{itemize}[leftmargin=*]
    \item \textbf{Estimated Time}: 2-3 minutes total
    \begin{itemize}
        \item Understanding problem: 20 seconds
        \item Computing trig ratios: 30 seconds
        \item Applying double angle formula: 40 seconds
        \item Recognizing pattern and applying cofunction: 30 seconds
        \item Verification: 20 seconds
        \item Marking answer: 10 seconds
    \end{itemize}
    
    \item \textbf{Time Checkpoints}: 
    \begin{itemize}
        \item At 1 min: Should have computed basic trig ratios
        \item At 2 min: Should have tried double angle or angle sum formula
        \item At 2.5 min: Should have answer or be testing choices
        \item At 3 min: MUST have answer marked
    \end{itemize}
    
    \item \textbf{Priority Level}: HIGH
    \begin{itemize}
        \item P10 is very accessible for prepared students
        \item Worth full effort before moving to harder problems
        \item Standard trigonometry - high success rate
    \end{itemize}
\end{itemize}

\subsection*{B. Risk Assessment}
\begin{itemize}[leftmargin=*]
    \item \textbf{Confidence Level}: 95\%+
    \begin{itemize}
        \item Clean algebraic solution
        \item Multiple verification methods confirm answer
        \item Answer matches natural form in choice (C)
        \item All other choices eliminated
    \end{itemize}
    
    \item \textbf{Abandonment Criteria}: 
    \begin{itemize}
        \item If after 3 minutes no progress $\rightarrow$ test answer choices systematically
        \item If stuck on algebra $\rightarrow$ try geometric approach or answer testing
        \item Should NOT abandon - this is accessible P10
    \end{itemize}
    
    \item \textbf{Flag for Review}: NO (very high confidence, multiple confirmations)
\end{itemize}

\subsection*{C. Competition Strategy Notes}
\begin{itemize}[leftmargin=*]
    \item \textbf{Pattern Recognition}: The appearance of $2\alpha$ in answer choice (C) should immediately suggest trying double angle formulas
    
    \item \textbf{Answer Choice Clues}: 
    \begin{itemize}
        \item (C) involves $\pi/2$ and $2\alpha$ - suggests complementary and double angle
        \item Other choices involve fractions or specific $\pi$ values - less natural
    \end{itemize}
    
    \item \textbf{Pythagorean Triple Recognition}: Both $3{-}4{-}5$ and $7{-}24{-}25$ are well-known Pythagorean triples - this guarantees exact rational trig values, making the problem algebraically clean
\end{itemize}
\end{competitionbox}

\begin{keyinsightbox}
\textbf{Key Insight}: The crucial observation is that $\cos(2\alpha) = \frac{7}{25} = \sin\beta$. This equality immediately suggests using the cofunction identity $\cos\theta = \sin(\pi/2 - \theta)$ to establish that $2\alpha$ and $\beta$ are complementary angles. The double angle formula transforms the problem from finding a relationship between two different triangles into recognizing a complementary angle relationship.
\end{keyinsightbox}

\begin{warningbox}
\textbf{Warning}: Don't confuse the smallest angle with the right angle! The smallest angle in a right triangle is always opposite the shortest side, not the right angle itself. Also, be careful with the double angle formula - use $\cos(2\alpha) = \cos^2\alpha - \sin^2\alpha$ rather than the $2\cos^2\alpha - 1$ form to avoid sign errors. Finally, when using $\sin\beta = \sin(\pi/2 - 2\alpha)$, verify that both angles are in the range where sine is injective (which they are, since both are acute angles).
\end{warningbox}

\begin{tipbox}
\textbf{Pro Tip}: For AMC trigonometry problems involving Pythagorean triples, always compute all three basic trig ratios ($\sin$, $\cos$, $\tan$) immediately. Then try standard identities: double angle, half angle, angle sum/difference. When you see answer choices involving $\pi/2$, think ``complementary angles'' and cofunction identities. For P9-16 problems, if the algebra gets messy after 2 minutes, pivot to testing answer choices - it's often faster and equally reliable.
\end{tipbox}

\section*{Final Answer}

\boxed{\textbf{(C) } \frac{\pi}{2} - 2\alpha}

\section*{Confidence Assessment}
\begin{itemize}[leftmargin=*]
    \item \textbf{Confidence Level}: 95-98\%
    \begin{itemize}
        \item Double angle formula yielded exact match: $\cos(2\alpha) = 7/25 = \sin\beta$
        \item Cofunction identity cleanly applied: $\sin\beta = \sin(\pi/2 - 2\alpha)$
        \item Alternative method (angle addition) confirms same answer
        \item Geometric construction method also yields same result
        \item All other answer choices can be eliminated numerically
    \end{itemize}
    
    \item \textbf{Verification Applied}: Level 3 - Targeted Spot Check
    \begin{itemize}
        \item Verified numerical equality: $\cos(2\alpha) = \sin\beta = 7/25$
        \item Confirmed angle sum: $\sin(2\alpha + \beta) = 1$ as expected for $\pi/2$
        \item Tested and eliminated all other answer choices
    \end{itemize}
    
    \item \textbf{Flag for Review}: No
    
    \item \textbf{Time Spent}: Approximately 2-3 minutes (optimal for P10 mid-band)
\end{itemize}

\end{document}

